% ECONOMICS_PROBLEM: {"id": "R23-505", "title": "Gentrification and incumbent welfare", "jel_code": "R23", "source_id": "OP-0465"}
% !TeX program = xelatex
\documentclass[11pt,a4paper]{article}
\usepackage[margin=23mm]{geometry}
\usepackage{fontspec}
\setmainfont{Latin Modern Roman}
\usepackage{amsmath,amssymb,mathtools}
\usepackage{xcolor,enumitem,booktabs,longtable,array,xurl,hyperref}
\definecolor{muted}{HTML}{53636C}
\hypersetup{colorlinks=true,urlcolor=blue,linkcolor=blue}
\setlength{\parindent}{0pt}
\setlength{\parskip}{5pt}
\newcommand{\R}{\mathbb R}
\newcommand{\E}{\mathbb E}
\newcommand{\N}{\mathbb N}
\newcommand{\ind}{\mathbf 1}
\newcommand{\Var}{\operatorname{Var}}
\newcommand{\Cov}{\operatorname{Cov}}
\newcommand{\supp}{\operatorname{supp}}
\DeclareMathOperator*{\argmin}{arg\,min}
\DeclareMathOperator*{\argmax}{arg\,max}
\newcommand{\entrymeta}[3]{{\sffamily\small\color{muted}\textbf{\texttt{#1}}\quad|\quad #2\par #3\par}}
\newcommand{\Problemref}[1]{\texttt{#1}}
\newcommand{\JELref}[2]{\texttt{#2} (JEL \texttt{#1})}
\begin{document}
\section*{Gentrification and incumbent welfare}
\textbf{Problem ID:} \texttt{R23-505}\quad \textbf{JEL:} \texttt{R23}\par
% BEGIN ECONOMICS STATEMENT
\label{problem:OP-0465}
{\sffamily\small\color{muted}Original subject group: Urban, housing and spatial economics\par}
\label{definitions:problem:30007747}
\textbf{Audit classification:} Published research agenda. Checked source identity and appropriate agenda/background support; expanded intervention, units, population, definedness and causal restrictions. Exact targets and variants are editorial.
\subsection*{Definitions and precise research target}
\textit{Editorial formalization.} Consider households living in a prespecified set of low-income urban neighborhoods immediately before a revitalization program is announced. Let \(a=1\) denote the defined investment package and \(a=0\) its absence. For baseline household \(i\), let \(V_i(a)\) be ten-year monetary-equivalent welfare combining consumption, housing quality, travel, safety, and valued neighborhood ties under the realized residential path. Define displacement \(D_i(a)\in\{0,1\}\) as a move caused by eviction or an unaffordable housing cost increase. The targets are
\[\tau_V=E[V_i(1)-V_i(0)],\qquad\tau_D=E[D_i(1)-D_i(0)].\]
Follow every baseline household after moving; repeated cross-sectional neighborhood averages do not identify these incumbent outcomes. Specify credible program variation, announcement effects, and counterfactual trends, and measure nearby spillovers and landlord capital gains separately. Distinguish voluntary moves from displacement using a prespecified classification with sensitivity analysis. Monetizing social ties requires explicit preference assumptions or bounds, rather than treating every move as a welfare loss.

The population is every baseline household in the selected neighborhoods before announcement, tracked at its later locations. Define \(\mathcal U_i^a\) on ten-year consumption, housing, safety and social-tie paths, including income/rent changes. Set \(V_i(a)\) as a baseline-money equivalent using a fixed feasible transfer schedule and strict monotonicity, or report utility/bounded outcomes separately when social ties cannot be monetized. Define \(D_i(a)\) from observable eviction or an announced rent-burden threshold followed by a move; a move's unobserved individual causal reason is not automatically measurable. Then \(\tau_D\) is the effect on this operational displacement indicator. Policy effect on all moves is a distinct estimand. Attrition and emigrants must be included or bounded; changing neighborhood composition cannot stand in for incumbent welfare. All expectations use a stated baseline population and probability law. Identification means every admissible data-generating model with the same observed-data law yields the same displayed target; otherwise report the identified set. Exact equations, horizons and assumptions are editorial, rather than source-given canonical conjectures.
\subsection*{Variants and assumption changes}
\begin{enumerate}
\item \textbf{Editorial tenure variant.} Estimate baseline homeowner/renter contrasts separately. Ownership gains and renter displacement cannot be inferred from neighborhood mean-income changes.
\item \textbf{Editorial spillover variant.} Include adjacent areas and displacement destinations using predeclared geography. Regional benefits differ from treated-area incumbent benefits.
\end{enumerate}
\subsection*{References and source audit}
\textbf{1.} Explicit agenda on mechanism weights, scaling mobility programs, and incumbent residents of revitalized areas. Locator: Journal of Economic Perspectives 35(4), 197–222; Discussion and Conclusion, pp.216–217. Full text or primary publication page inspected. URL: \url{https://lkatz.scholars.harvard.edu/sites/g/files/omnuum5961/files/lkatz/files/chyn_katz_jep_2021_all.pdf}.
\begin{enumerate}\raggedright
\item\hypertarget{definitions:source:30007747:1}{} Eric Chyn; Lawrence F. Katz. 2021. \emph{Neighborhoods Matter: Assessing the Evidence for Place Effects}. Journal of Economic Perspectives 35(4), 197–222; Discussion and Conclusion, pp.216–217.
\url{https://lkatz.scholars.harvard.edu/sites/g/files/omnuum5961/files/lkatz/files/chyn_katz_jep_2021_all.pdf}.
DOI: \url{https://doi.org/10.1257/jep.35.4.197}.
\end{enumerate}

\subsection*{Common conventions: Source mathematical conventions}

These conventions apply to each entry unless explicitly replaced.

\textbf{Models and identification.} A primitive model \(m\) specifies all probability laws, feasible choices, information, equilibrium and measurement rules used by its target. The admissible class \(\mathcal M\) is fixed independently of outcomes. If \(P_O(m)\) is its observable probability law and \(\tau(m)\) the target, its identified set at observable law \(P\) is
\[
 \mathcal I_\tau(P)=\{\tau(m):m\in\mathcal M,\ P_O(m)=P\}.
\]
Point identification means that this set is a singleton. An empty set signals incompatibility of data and assumptions. Coordinate bounds need not describe the joint set. Bounds are sharp only if every retained target is attainable or arbitrarily approachable by compatible models. Finite bounds require explicit restrictions on unobserved quantities; unrestricted confounding, hidden wealth, extrapolation or arbitrary equilibrium selection can yield uninformative sets.

\textbf{Causal interventions.} A potential outcome \(Y_i(p)\) is the outcome under a complete policy regime \(p\), including timing, financing, information and market spillovers. The target population and units are fixed before treatment. With interference, use the complete assignment vector or a justified exposure mapping; individual no-interference notation alone cannot identify an economy-wide intervention. A difference of mean potential outcomes does not require the cross-world joint distribution. Conditioning on post-treatment migration, survival, enrollment or employment changes the estimand and needs separate justification.

\textbf{Statistical statements.} An estimator, confidence set, forecast or test specifies a sampling law, model class and loss/coverage criterion. Uniform \((1-\alpha)\)-coverage means \(\inf_{m\in\mathcal M}\Pr_m\{\tau(m)\in C_n\}\geq1-\alpha\), or an explicitly asymptotic version. The whole parameter space trivially has coverage, so informativeness requires a stated width, risk or power objective. A forecasting improvement is relative to a named benchmark, information set and evaluation population.

\textbf{Equilibrium and optimization.} An equilibrium comprises feasible plans, optimality, beliefs where needed, budget constraints and all market/resource clearing. The primitives or a stated benchmark define each correspondence \(\mathcal E(p;m)\). Nonemptiness is assumed or proved, never inferred just from compact policy sets. A selected-equilibrium target uses a declared selection rule; a robust target quantifies over all permitted equilibria. Use suprema or infima unless attainment is established. An \(\varepsilon\)-optimal policy is feasible and within \(\varepsilon\) of the finite optimum value. Report an empty feasible set separately.

\textbf{Welfare and accounting.} Normative utility, interpersonal normalization, social weights, numeraire, discounting and the use of public receipts are primitives. Behavioral utility need not coincide with the welfare criterion. Household welfare includes ownership income and rents once; adding profits again to compensating variation generally double-counts them. A monetary equivalent is defined by an explicit transfer and price convention, with monotonicity and range conditions. Average effects, mechanism contributions, transfers and welfare losses are different targets. Infinite sums require convergence and dynamic feasible plans require terminal or no-Ponzi conditions.

\textbf{Source versions.} A working-paper date, revision date and journal publication year are distinguished. A DOI for a journal version is not automatically the DOI of an earlier manuscript. Locators refer to the stated version and printed pages where specified. A metadata-only check does not verify a claim about a full-text conclusion. Every entry's source subsection narrows the provenance claim accordingly.
% END ECONOMICS STATEMENT
\end{document}
